\uuid{njSj}
\exo7id{1833}
\titre{exo7 1833}
\auteur{ridde}
\organisation{exo7}
\datecreate{1999-11-01}
\isIndication{false}
\isCorrection{false}
\chapitre{Extremum, extremum lié}
\sousChapitre{Extremum, extremum lié}
\cpge{mpsi-mp2i.derivabilite.b, pcsi.derivabilite.b, ptsi.derivabilite.b, mpsi-mp2i.analyse-asymptotique.b, pcsi.analyse-asymptotique.b, ptsi.analyse-asymptotique.b, mpsi-mp2i.fonctions-deux-variables.b, pcsi.fonctions-deux-variables.b, ptsi.fonctions-deux-variables.b, mpsi-mp2i.fonctions-deux-variables.d, pcsi.fonctions-deux-variables.d, ptsi.fonctions-deux-variables.d, bcpst1.developpements-limites.a, bcpst1.fonctions-deux-variables.c, mp-mpi.calcul-differentiel-optimisation.f}
\module{Calcul différentiel}
\niveau{L3}
\difficulte{}

\contenu{
\question{
Soit ${f}: {\Rr^{2}} \to {\Rr}$, ${ (x, y)}\mapsto {xe^{y} + ye^{x}}$.\\
 Montrer que $ (-1, -1)$ est le seul extremum possible. A l'aide d'un
 d\'eveloppement limit\'e de $\varphi (h) = f (-1 + h, -1 + h)$ et de
 $\psi (h) = f (-1 + h, -1-h)$, montrer que $f$ n'a pas d'extremum.
}
}
